\section{Transfer to Frozen DNABERT-2 Promoter Classification}
\label{sec:dnabert}

We next test whether the same behavior-first causal workflow transfers beyond SpliceBERT to a second
genomic foundation model. We analyze a frozen DNABERT-2 encoder on promoter classification, using
masked-mean pooled hidden states and a separately trained linear probe. This readout measures where
promoter information is linearly accessible in the frozen representation; it is not a native pretrained
DNABERT-2 output head. Candidate components are selected on validation data and evaluated on held-out
test examples.

\subsection{Frozen Representation and Corruption Localization}
\label{sec:dnabert-representation}

A corrected layer-wise probe sweep selects Layer~2 as the strongest frozen readout endpoint. On the
held-out test split, the Layer-2 probe reaches approximately \(0.810\) balanced accuracy and \(0.892\)
ROC--AUC, outperforming the strongest tested k-mer baseline. An earlier analysis selected Layer~5, but
that endpoint was superseded after the frozen pipeline was corrected and rerun.

\begin{figure}[t]
    \centering
    \includegraphics[width=0.94\linewidth]{images/figure1_frozen_layer_probe_validation.png}
    \caption{\textbf{Frozen DNABERT-2 layer-wise promoter probing.}
    Validation balanced accuracy and ROC--AUC peak at Layer~2, which is fixed as the downstream
    readout for the subsequent frozen-model analysis. The selected Layer-2 probe reaches
    \(0.810\) balanced accuracy and \(0.892\) ROC--AUC on the held-out test split.}
    \label{fig:dnabert-layer-probe}
\end{figure}

Sequence corruption provides a complementary localization assay. Mononucleotide shuffling produces a
substantially larger balanced-accuracy drop than dinucleotide-preserving shuffling, and the largest
mononucleotide damage occurs at Layer~3 on both validation and held-out test trajectories. We therefore
distinguish the strongest \emph{readout} layer (Layer~2) from the strongest corruption-sensitive
\emph{processing} stage (Layer~3); the latter motivates component-level intervention rather than replacing
the locked Layer-2 endpoint.

\begin{figure}[t]
    \centering
    \includegraphics[width=0.94\linewidth]{images/figure3_test_corruption_trajectory_drops.png}
    \caption{\textbf{Held-out DNABERT-2 corruption sensitivity.}
    Balanced-accuracy damage under mononucleotide and dinucleotide-preserving corruption across layers.
    Mononucleotide damage peaks at Layer~3, motivating targeted causal analysis of early components while
    leaving Layer~2 fixed as the predictive readout.}
    \label{fig:dnabert-corruption}
\end{figure}

\subsection{Component-Level Causal Localization}
\label{sec:dnabert-causal}

We next intervene on attention heads and structured groups in the early frozen network. Individual
head ablations reveal heterogeneous effects, with several promoter-supporting heads concentrated in the
early layers. These scans are used to generate candidates rather than to define a circuit by ranking
alone.

\begin{figure}[t]
    \centering
    \includegraphics[width=0.94\linewidth]{images/figure4_layer3_head_heatmap.png}
    \caption{\textbf{Frozen DNABERT-2 head-ablation map.}
    Promoter mean damage under head ablation across the tested early layers. More negative values indicate
    components whose removal reduces promoter evidence. Marked heads are candidate components taken
    forward to structured interventions.}
    \label{fig:dnabert-head-map}
\end{figure}

Grouped ablation provides a stronger test of collective necessity. On the held-out test cohort, both the
expanded pathway and the primary pathway damage promoter evidence in the large majority of promoter
examples. This supports reproducible collective dependence on the selected component groups, but does
not imply that the computation is sparse or that any single group is complete.

\begin{figure}[t]
    \centering
    \includegraphics[width=0.90\linewidth]{images/figure6_test_grouped_ablation_fractions.png}
    \caption{\textbf{Held-out grouped-ablation consistency.}
    Fraction of promoter examples damaged by the validation-selected expanded and primary Layer-3
    pathways. Both interventions affect most held-out promoters, supporting collective necessity of the
    selected groups.}
    \label{fig:dnabert-grouped-ablation}
\end{figure}

\subsection{Donor Transfer and Information Specificity}
\label{sec:dnabert-donor}

Ablation shows necessity but not whether a component carries recoverable promoter-relevant information.
We therefore use donor-transfer interventions in which candidate internal activations are transplanted
under matched and mismatched donor conditions. On held-out test examples, matched-promoter and same-class
promoter donors yield positive mean recovery, whereas opposite-class non-promoter donors do not. This
class-structured contrast is evidence that the tested pathway carries recoverable information related to
the promoter decision rather than only generic activation magnitude.

\begin{figure}[t]
    \centering
    \includegraphics[width=0.90\linewidth]{images/figure8_test_donor_matrix.png}
    \caption{\textbf{Held-out Layer-3 donor-transfer matrix.}
    Matched-promoter and same-class promoter donors produce positive recovery, whereas opposite-class
    non-promoter donors do not. Zero activation provides a magnitude control and self-corrupt is the
    reference baseline. The contrast supports class-structured recoverable information in the tested
    pathway.}
    \label{fig:dnabert-donor-matrix}
\end{figure}

\subsection{Biological Perturbation and Evidence Boundary}
\label{sec:dnabert-tata}

We separately test whether the frozen promoter representation reduces to sensitivity to canonical
TATA-like promoter motifs. The current validation cohort contains 71 eligible promoters matching
\texttt{TATA[AT][AT]}; 48 contain the canonical \texttt{TATAAA} pattern. The mean excess effect is
\(0.074\) across all eligible TATA-like cases and \(0.145\) within the \texttt{TATAAA} subgroup, but the
effects are heterogeneous and not statistically conclusive under the current assay. We therefore do not
interpret the frozen DNABERT-2 mechanism as a TATA-box detector.

\begin{table}[t]
    \centering
    \caption{\textbf{Exploratory frozen DNABERT-2 TATA-like perturbation results.}
    These validation-stage results are treated as a biological hypothesis test, not as a confirmed
    semantic interpretation of the causal components.}
    \label{tab:dnabert-tata}
    \small
    \begin{tabular}{lrrrl}
        \toprule
        Cohort & \(n\) & Mean excess & Fraction \(>0\) & Interpretation \\
        \midrule
        All TATA-like & 71 & 0.074 & 0.521 & Not conclusive \\
        TATAAA & 48 & 0.145 & 0.563 & Suggestive, not conclusive \\
        TATATA & 14 & 0.027 & -- & Not conclusive \\
        TATAAT & 4 & -0.549 & -- & Too small \\
        TATATT & 5 & 0.029 & -- & Too small \\
        \bottomrule
    \end{tabular}
\end{table}

Taken together, the DNABERT-2 extension establishes a second-model transfer of the experimental
workflow: promoter information is predictively accessible in the frozen model, corruption identifies an
early sensitive stage, structured ablations test collective necessity, and donor interventions test
recoverable information. The present evidence does not establish a fully mapped promoter-recognition
circuit or a biologically identified TATA-specific mechanism. A matched fine-tuned DNABERT-2 analysis is
being developed separately to test whether supervised training preserves, concentrates, or reorganizes
this frozen computation.
